% Dokumen pemanasan cache Tectonic.
%
% Dipakai saat BUILD IMAGE (lihat docker/server.Dockerfile) untuk mengunduh
% paket-paket yang paling sering dipakai, sehingga compile pertama pengguna
% tidak perlu menunggu unduhan.
%
% Sengaja memuat paket yang dipakai template bawaan aplikasi (amsmath,
% graphicx, hyperref) plus paket umum lain. Tanpa ini, cache bawaan image hanya
% berisi article.cls sehingga compile pertama memakan ~30 detik.
\documentclass[11pt]{article}

% --- paket template bawaan aplikasi ---
\usepackage[utf8]{inputenc}
\usepackage[T1]{fontenc}
\usepackage{amsmath, amssymb}
\usepackage{graphicx}
\usepackage{hyperref}

% --- paket umum lain ---
\usepackage{geometry}
\usepackage{xcolor}
\usepackage{booktabs}
\usepackage{listings}
\usepackage{caption}
\usepackage{enumitem}
\usepackage{float}
\usepackage{multirow}
\usepackage{array}
\usepackage{url}
\usepackage{verbatim}
\usepackage{fancyhdr}
\usepackage{titlesec}
\usepackage{parskip}
\usepackage[english]{babel}

\begin{document}

\title{Pemanasan cache}
\author{LaTeX Web Editor}
\date{\today}
\maketitle

\begin{abstract}
Ringkasan singkat.
\end{abstract}

\tableofcontents

\section{Pendahuluan}
Teks biasa dengan rumus inline $E = mc^2$ dan rujukan silang
\ref{eq:integral}.

\section{Metode}
\begin{equation}
  \int_0^1 x^2 \, dx = \frac{1}{3}
  \label{eq:integral}
\end{equation}

\begin{align}
  a &= b + c \\
  d &= e - f
\end{align}

\begin{itemize}
  \item Butir pertama
  \item Butir kedua
\end{itemize}

\begin{enumerate}
  \item Nomor satu
  \item Nomor dua
\end{enumerate}

\begin{table}[H]
  \centering
  \begin{tabular}{lcr}
    \toprule
    Kiri & Tengah & Kanan \\
    \midrule
    a & b & c \\
    \bottomrule
  \end{tabular}
  \caption{Contoh tabel}
\end{table}

\begin{verbatim}
contoh verbatim
\end{verbatim}

\begin{lstlisting}[language=Python]
def halo():
    return "hai"
\end{lstlisting}

\section{Kesimpulan}
Selesai. Lihat \url{https://example.com}.

\end{document}
